\title{Group Sequence Policy Optimization}

\begin{abstract}
We present Group Sequence Policy Optimization (GSPO), a reinforcement learning algorithm designed for training large language models with enhanced stability, efficiency, and effectiveness. In contrast to prior methods that rely on importance ratios at the token level, GSPO instead calculates the importance ratio at the sequence level, applying clipping, reward assignment, and optimization processes to entire sequences. Through comprehensive evaluation, we show that GSPO outperforms the GRPO algorithm in both training efficiency and end performance, offers marked improvements in stabilizing Mixture-of-Experts (MoE) reinforcement learning, and opens opportunities to streamline RL system architecture. These advantages have played a pivotal role in driving significant progress in the development of the most recent Qwen3 models.
\end{abstract}

\section{Introduction}
\label{sec:intro}

Reinforcement learning (RL) has become a central approach in advancing the scalability of language models \citep{o1,dpsk-r1,qwq32b,qwen3}. Utilizing large-scale RL enables these models to address complex tasks—including those in advanced mathematics and programming—by engaging in more extensive and intricate reasoning.

In order to effectively leverage increased computational resources for scaling RL, it is essential to ensure training processes remain both stable and resilient. Nonetheless, leading RL techniques such as GRPO \citep{grpo} encounter significant stability challenges when applied to the training of very large language models, frequently culminating in drastic and unrecoverable failures of the model \citep{qwen3, minimax-m1}. Such persistent instability impedes advancements in enhancing language model performance via extended RL-based optimization.

This study reveals that GRPO's instability originates from a basic flaw in its algorithmic approach—namely, the improper use and breakdown of importance sampling weights. This flaw gives rise to substantial training variance, which becomes more pronounced as response lengthens and is intensified by the clipping procedure, eventually leading to model failure.

In response to these fundamental challenges, we introduce \textbf{Group Sequence Policy Optimization (GSPO)}, an RL method designed specifically for optimizing large language models. GSPO's primary advancement is its theoretically motivated importance ratio, which is formulated through sequence likelihood in line with established importance sampling principles \citep{click}. Furthermore, GSPO determines advantages by normalizing rewards across several responses for each query, thereby promoting consistency between sequence-level rewards and the learning objective.

Our experimental results highlight the marked advantages of GSPO compared to GRPO, particularly in terms of stability during training, computational efficiency, and overall effectiveness. Importantly, GSPO inherently addresses the instability issues commonly encountered in reinforcement learning when applied to large Mixture-of-Experts (MoE) architectures, obviating the requirement for intricate stabilization mechanisms and offering the promise of streamlined RL frameworks. These benefits have played a pivotal role in driving the substantial gains observed in the latest iterations of Qwen3 models. We anticipate that GSPO will serve as a solid and scalable cornerstone for future progress in the training of large-scale language models with reinforcement learning.

\section{Preliminaries}

\paragraph{Notation}
Throughout this work, we represent an autoregressive language model, parameterized by $\theta$, as a policy $\pi_\theta$. The variable $x$ indicates a query, while $\mathcal{D}$ refers to the set of all queries. For a particular response $y$ corresponding to the query $x$, the probability of $y$ under the policy $\pi_\theta$ is written as $\pi_\theta (y | x)=\prod_{t=1}^{|y|} \pi_\theta (y_t | x, y_{<t} )$, where $|y|$ specifies the total number of tokens in $y$. Additionally, a verifier function $r$ assigns a reward $r(x, y) \in [0, 1]$ to each query-response pair $(x, y)$.

\paragraph{Proximal Policy Optimization (PPO)} Proximal Policy Optimization (PPO) \citep{ppo} utilizes trajectories collected under the previous policy $\pi_{\theta_\text{old}}$ and restricts policy updates to remain close to this original policy by means of a clipping technique. More precisely, PPO optimizes the policy according to the objective function below (the KL regularization component is excluded in subsequent discussion for clarity, since it is not central to our analysis):

\begin{align}
\mathcal{J}_\text{PPO}(\theta) = \mathbb{E}_{ x \sim \mathcal{D},\, y \sim \pi_{\theta_\text{old}}( \cdot | x) }
\left[ \frac{1}{|y|} \sum_{t=1}^{|y|} 
\min \left( w_{t}(\theta) \widehat{A}_{t},  \, \mathrm{clip} \left( w_{t}(\theta), 1 - {\varepsilon}, 1 + {\varepsilon}\right) \widehat{A}_{t} \right)
\right],
\end{align}

Here, the token $y_t$'s importance ratio is expressed as $w_{t}(\theta) = \frac{ \pi_{\theta} (y_{t} | x, y_{<t}) }{ \pi_{\theta_\text{old}} (y_{t} | x,y_{<t})}$, where $\widehat{A}_{t}$, representing the advantage of $y_t$, is obtained via a separate value model, and $\varepsilon$ denotes the limit for clipping the importance ratios.

A principal practical limitation of PPO arises from its significant dependence on the value model. In particular, the value model is typically as large as the policy model, thereby imposing substantial demands on memory and computation. Moreover, the algorithm’s overall performance is highly sensitive to the accuracy of its value estimates. Not only is training a dependable value model a difficult endeavor by itself, but extending its reliability to accommodate longer outputs and more intricate tasks proves to be even more formidable.

\paragraph{Group Relative Policy Optimization (GRPO)} GRPO \citep{grpo} eliminates reliance on a value model by assessing the relative advantage of individual responses compared to others within a set addressing the same prompt. In particular, GRPO maximizes the following objective:

\begin{align}
\label{equ:grpo}
\mathcal{J}_\text{GRPO}(\theta) = \mathbb{E}_{ x \sim \mathcal{D},\, \{y_i\}_{i=1}^G \sim \pi_{\theta_\text{old}}( \cdot | x) }
\left[ \frac{1}{G} \sum_{i=1}^{G} \frac{1}{|y_i|} \sum_{t=1}^{|y_i|} 
\min \left( w_{i,t}(\theta) \widehat{A}_{i,t},  \, \mathrm{clip} \left( w_{i,t}(\theta), 1 - {\varepsilon}, 1 + {\varepsilon}\right) \widehat{A}_{i,t} \right)
\right],
\end{align}

Here, $G$ denotes the number of responses produced per input query $x$ (representing the group size), while $w_{i,t}(\theta)$ indicates the importance ratio and $\widehat{A}_{i,t}$ corresponds to the advantage computed for token $y_{i,t}$ as follows:

\begin{align}
    w_{i,t}(\theta)=\frac{ \pi_{\theta} (y_{i,t} | x, y_{i,<t}) }{ \pi_{\theta_\text{old}} (y_{i,t} | x,y_{i,<t})},\quad\ 
    \widehat{A}_{i,t} = \widehat{A}_{i} = \frac{r(x, y_i) - \mathrm{mean} \left( \{ r(x, y_i) \}_{i=1}^G \right) }{ \mathrm{std} \left( \{ r(x, y_i) \}_{i=1}^G \right) },
\end{align}

respectively, wherein every token within $y_i$ possesses the identical advantage value denoted by $\widehat{A}_{i}$.

\section{Motivation}
\label{sec:motivation}

As models increase in size and sparsity—such as those found in Mixture-of-Experts architectures—and as the length of responses grows, reinforcement learning demands large rollout batch sizes to fully leverage computational resources. To enhance the efficiency with which samples are used, it is common to divide these large rollout batches into several smaller mini-batches for performing gradient updates. This approach, however, results in an off-policy learning context, since the sampled responses $y$ originate from a previous policy, $\pi_{\theta_\text{old}}$, instead of the current policy $\pi_\theta$ under optimization. Consequently, this shift underscores the importance of incorporating clipping strategies, as in PPO and GRPO, to exclude samples that are excessively ``off-policy'' from the gradient computation process.

Although approaches such as clipping seek to control the off-policy discrepancy, we observe a deeper flaw in GRPO: \textit{the formulation of its objective is fundamentally ill-posed}. This challenge is especially pronounced when employing large-scale models for tasks requiring extended outputs, which can result in severe degradation of model performance. The ill-posedness of the GRPO objective arises from the incorrect use of importance sampling weights. Importance sampling is intended to approximate the expectation of a function $f$ with respect to a target distribution $\pi_\text{tar}$ by appropriately weighting samples that originate from a behavior distribution $\pi_\text{beh}$:

\begin{align}
\mathbb{E}_{z \sim \pi_\text{tar}} \left[ f(z) \right] = \mathbb{E}_{z \sim \pi_\text{beh}} \left[ \frac{ \pi_\text{tar} (z) }{ \pi_\text{beh} (z)} \, f(z) \right].
\end{align}

A key aspect of this approach is the necessity to average across many samples ($N \gg 1$) drawn from the behavior distribution $\pi_\text{beh}$ in order for the importance weight $\frac{ \pi_\text{tar} (z) }{ \pi_\text{beh} (z)}$ to adequately address differences between the distributions.

On the other hand, GRPO incorporates the importance weight $\frac{ \pi_{\theta} (y_{i,t} | x, y_{i,<t}) }{ \pi_{\theta_\text{old}} (y_{i,t} | x,y_{i,<t})}$ at each step $t$ of the token sequence. However, because this weighting relies on a single sampled token $y_{i,t}$ from the corresponding next-token probability distribution $\pi_{\theta_\text{old}}( \cdot | x, y_{i, <t})$, it does not effectively correct for distributional mismatch as intended. Rather, it injects considerable high-variance noise into the gradient estimates during training, with the effect becoming more severe for longer sequences and further amplified by the use of a clipping mechanism. Through empirical analysis, we have found that such issues may cause the model to collapse—a failure mode that often cannot be remedied. After collapse, training cannot be successfully restarted, even by reverting to earlier model checkpoints, making extensive adjustments to hyperparameters (such as clipping bounds), increasing the length of generated sequences, or altering RL query strategies.

The preceding analysis highlights a basic flaw inherent in GRPO's structure. Specifically, the ineffectiveness of the token-level importance weights reveals an underlying rule: \textbf{the granularity of the optimization target ought to correspond with that of the reward}. Given that rewards are allocated at the sequence level, conducting off-policy corrections for individual tokens introduces complications. Consequently, we are motivated to abandon the token-level optimization framework and instead consider employing importance weighting and optimizing at the \textit{sequence level}.

\section{Algorithm}

\subsection{GSPO: Group Sequence Policy Optimization}

Although the use of the token-level importance weight $\frac{ \pi_{\theta} (y_{i,t} | x, y_{i,<t}) }{ \pi_{\theta_\text{old}} (y_{i,t} | x,y_{i,<t})}$ poses challenges in GRPO, our analysis indicates that, for language generation tasks, the \textit{sequence-level} importance weight $\frac{ \pi_{\theta} (y | x) }{ \pi_{\theta_\text{old}} (y | x)}$ possesses clear theoretical significance. Specifically, it quantifies the extent to which a response $y$—sampled according to $\pi_{\theta_\text{old}} (\cdot | x)$—differs from what is expected under $\pi_{\theta} (\cdot | x)$. This notion is consistent with the sequence-level reward framework and provides a meaningful criterion for applying the clipping mechanism.

Motivated by this simple insight, we introduce the \textbf{Group Sequence Policy Optimization (GSPO)} algorithm. The central idea of GSPO is to utilize the following objective that optimizes at the sequence level:

\begin{align}
\mathcal{J}_\text{GSPO} (\theta) =
\mathbb{E}_{ x \sim \mathcal{D},\, \{y_i\}_{i=1}^G \sim \pi_{\theta_\text{old}}( \cdot | x) }
\left[ 
\frac{1}{G} \sum_{i=1}^{G}
\min \left( s_{i}(\theta)  \widehat{A}_{i},  \, \mathrm{clip} \left( s_{i}(\theta), 1 - {\varepsilon}, 1 + {\varepsilon} \right) \widehat{A}_{i} \right) 
\right],
\label{equ:gspo}
\end{align}

in which we utilize the group-level advantage estimation:

\begin{align}
\widehat{A}_{i} = \frac{r(x, y_i) - \mathrm{mean} \left( \{ r(x, y_i) \}_{i=1}^G \right) }{ \mathrm{std} \left( \{ r(x, y_i) \}_{i=1}^G \right) },
\end{align}

and introduce the importance ratio $s_{i}(\theta)$, which is formulated using sequence likelihood as in \citet{click}:

\begin{align}
s_{i}(\theta) = \left( \frac{ \pi_{\theta} (y_i | x) }{ \pi_{\theta_\text{old}} (y_i | x)} \right)^{\frac{1}{|y_i|}}
=
\exp \left( \frac{1}{|y_i|} \sum_{t=1}^{|y_i|} \log \frac{ \pi_{\theta} (y_{i,t} | x, y_{i,<t}) }{ \pi_{\theta_\text{old}} (y_{i,t} | x,y_{i,<t})} \right).
\end{align}

Consequently, GSPO performs clipping at the response level, rather than at the level of individual tokens, thereby filtering out excessively ``off-policy'' samples from the gradient estimation process. This approach aligns with both the reward calculation and optimization being conducted at the sequence level. Additionally, we implement length normalization in $s_{i}(\theta)$, aiming to decrease variance and confine $s_{i}(\theta)$ within a consistent numerical interval. Without this normalization, substantial changes in token likelihoods could cause significant swings in the importance ratio calculated across entire sequences, and the clipping intervals required for responses of varying lengths would not be uniform. It is important to emphasize that the scales of the clipping ranges in GSPO differ significantly—often by orders of magnitude—from those in earlier algorithms such as GRPO, primarily owing to the different formulations of importance ratios.

\subsection{Gradient Analysis}
\label{subsec:gradient}

The gradient of the GSPO objective may be computed in the following manner (with clipping disregarded for conciseness):

\begin{align}
\nabla_{\theta} \mathcal{J}_\text{GSPO} (\theta)
=&\ 
\nabla_{\theta} \mathbb{E}_{ x \sim \mathcal{D},\, \{y_i\}_{i=1}^G \sim \pi_{\theta_\text{old}}( \cdot | x) }
\left[ 
\frac{1}{G} \sum_{i=1}^{G} s_{i}(\theta) \widehat{A}_{i}
\right] \\
= &\ 
\mathbb{E}_{ x \sim \mathcal{D},\, \{y_i\}_{i=1}^G \sim \pi_{\theta_\text{old}}( \cdot | x) }
\left[ 
\frac{1}{G} \sum_{i=1}^{G}
s_{i}(\theta) \widehat{A}_{i}
\cdot \nabla_{\theta} \log s_{i}(\theta)
\right] \\
=&\ 
\mathbb{E}_{ x \sim \mathcal{D},\, \{y_i\}_{i=1}^G \sim \pi_{\theta_\text{old}}( \cdot | x) }
\left[ 
\frac{1}{G} \sum_{i=1}^{G} 
\left( \frac{ \pi_{\theta} (y_i | x) }{ \pi_{\theta_\text{old}} (y_i | x)} \right)^{\frac{1}{|y_i|}} \widehat{A}_{i} 
\cdot \frac{1}{|y_i|} \sum_{t=1}^{|y_i|} \nabla_{\theta} \log \pi_{\theta} (y_{i,t} | x, y_{i,<t}) 
\right].
\label{equ:gspo_grad}
\end{align}

By way of comparison, the gradient of the GRPO objective can be written as (observe that $\widehat{A}_{i,t} = \widehat{A}_{i}$):

\begin{align}
\nabla_{\theta} \mathcal{J}_\text{GRPO} (\theta) 
=&\ 
\nabla_{\theta} \mathbb{E}_{ x \sim \mathcal{D},\, \{y_i\}_{i=1}^G \sim \pi_{\theta_\text{old}}( \cdot | x) }
\left[ \frac{1}{G} \sum_{i=1}^{G} \frac{1}{|y_i|} \sum_{t=1}^{|y_i|} 
w_{i,t}(\theta) \widehat{A}_{i,t}
\right] \\
=&\
\mathbb{E}_{ x \sim \mathcal{D},\, \{y_i\}_{i=1}^G \sim \pi_{\theta_\text{old}}( \cdot | x) }
\left[ \frac{1}{G} \sum_{i=1}^{G} \widehat{A}_{i} 
\cdot \frac{1}{|y_i|} \sum_{t=1}^{|y_i|} 
\frac{ \pi_{\theta} (y_{i,t} | x, y_{i,<t}) }{ \pi_{\theta_\text{old}} (y_{i,t} | x,y_{i,<t})} 
\nabla_{\theta} \log \pi_{\theta} (y_{i,t} | x, y_{i,<t})  
\right].
\end{align}

The principal difference separating GSPO from GRPO concerns \textit{the approach each method takes to weighting the gradients of token log likelihoods}. In the case of GRPO, individual tokens are multiplied by an ``importance weight'' factor, expressed as $\frac{ \pi_{\theta} (y_{i,t} | x, y_{i,<t}) }{ \pi_{\theta_\text{old}} (y_{i,t} | x,y_{i,<t})}$. These weights are inherently uneven—falling within the interval $(0, 1+\varepsilon]$ when $\widehat{A}_{i} > 0$, and within $[1-\varepsilon, +\infty)$ for $\widehat{A}_{i} < 0$—and their non-uniformity is significant, potentially accumulating and causing unpredictable dynamics over the course of training. In comparison, GSPO assigns uniform weights to all tokens within a response, thereby removing the source of instability present in GRPO.

\subsection{GSPO-token: A Token-level Objective Variant}

In situations such as multi-turn RL, it can be beneficial to adjust the advantage at a granularity finer than the entire sequence. Therefore, we propose a token-level version of GSPO, referred to as \textbf{GSPO-token}, which facilitates advantage tailoring on a per-token basis:

\begin{align}
\mathcal{J}_\text{GSPO-token}(\theta)
=
\mathbb{E}_{ x \sim \mathcal{D},\, \{y_i\}_{i=1}^G \sim \pi_{\theta_\text{old}}( \cdot | x) }
\left[ 
\frac{1}{G} \sum_{i=1}^{G} \frac{1}{|y_i|} \sum_{t=1}^{|y_i|} 
\min \left( s_{i,t}(\theta) \widehat{A}_{i,t},  \, \mathrm{clip} \left( s_{i,t}(\theta), 1 - {\varepsilon}, 1 + {\varepsilon} \right) \widehat{A}_{i,t} \right) 
\right],
\label{equ:gspo-token}
\end{align}

in which

\begin{align}
s_{i,t}(\theta) 
= \mathrm{sg} \left[ s_{i}(\theta) \right]  \cdot \frac{ \pi_{\theta} (y_{i,t} | x, y_{i,<t}) }{ \mathrm{sg} \left[ \pi_{\theta} (y_{i,t} | x, y_{i,<t}) \right] },
\end{align}

and $\mathrm{sg}[\cdot]$ refers to extracting the numerical value while preventing any gradient flow, akin to the \texttt{detach} function in PyTorch. The gradient computation for GSPO-token can be expressed as follows:

\begin{align}
\nabla_{\theta} \mathcal{J}_\text{GSPO-token}(\theta)
=&\ 
\nabla_{\theta} \mathbb{E}_{ x \sim \mathcal{D},\, \{y_i\}_{i=1}^G \sim \pi_{\theta_\text{old}}( \cdot | x) }
\left[ 
\frac{1}{G} \sum_{i=1}^{G}
\frac{1}{|y_i|} \sum_{t=1}^{|y_i|} 
s_{i,t}(\theta) \widehat{A}_{i,t}
\right] \\
=&\ 
\mathbb{E}_{ x \sim \mathcal{D},\, \{y_i\}_{i=1}^G \sim \pi_{\theta_\text{old}}( \cdot | x) }
\left[ 
\frac{1}{G} \sum_{i=1}^{G} s_i (\theta) \cdot \frac{1}{|y_i|} \sum_{t=1}^{|y_i|} 
\widehat{A}_{i,t} \frac{ \nabla_{\theta} \pi_{\theta} (y_{i,t} | x, y_{i,<t}) }{ \pi_{\theta} (y_{i,t} | x, y_{i,<t}) }
\right] \\
=&\ 
\mathbb{E}_{ x \sim \mathcal{D},\, \{y_i\}_{i=1}^G \sim \pi_{\theta_\text{old}}( \cdot | x) }
\left[ 
\frac{1}{G} \sum_{i=1}^{G}  \left( \frac{ \pi_{\theta} (y_i | x) }{ \pi_{\theta_\text{old}} (y_i | x)} \right)^{\frac{1}{|y_i|}}
\cdot \frac{1}{|y_i|} \sum_{t=1}^{|y_i|}
\widehat{A}_{i,t} \nabla_{\theta} \log \pi_{\theta} (y_{i,t} | x, y_{i,<t})  
\right].
\label{equ:gspo-token_grad}
\end{align}

Observe that the ratio $\frac{ \pi_{\theta} (y_{i,t} | x, y_{i,<t}) }{ \mathrm{sg} \left[ \pi_{\theta} (y_{i,t} | x, y_{i,<t}) \right] }$ always evaluates to 1, which implies $s_{i,t}(\theta)$ and $s_{i}(\theta)$ are equivalent numerically. By examining Equation~\eqref{equ:gspo} alongside~\eqref{equ:gspo-token}, as well as Equation~\eqref{equ:gspo_grad} and~\eqref{equ:gspo-token_grad}, it becomes evident that, when all tokens in the sequence $y_i$ share a common advantage value (i.e., setting $\widehat{A}_{i,t} = \widehat{A}_{i}$ for all $t$), GSPO-token is mathematically the same as GSPO in terms of optimization target, clipping mechanism, and the form of the theoretical gradient. Nevertheless, GSPO-token provides enhanced adaptability by enabling token-level specification of the advantages.

\section{Experiments and Discussion}

\subsection{Empirical Results}

We conduct experiments using a cold-start model that has been fine-tuned from Qwen3-30B-A3B-Base. For evaluation, we present both the training reward trajectories and the performance trends of the model on several benchmarks, including AIME'24 (reporting the mean Pass@1 across 32 samples), LiveCodeBench (202410-202502, with mean Pass@1 across 8 samples), and CodeForces (measured by Elo Rating). In the reinforcement learning training phase, each rollout batch is subdivided into four smaller mini-batches, which are used to compute gradient updates. For GSPO, the left and right clipping intervals in Equation~\eqref{equ:gspo} are configured to 3e-4 and 4e-4, respectively. Our experiments use GRPO as the reference baseline, where the left and right clipping intervals in Equation~\eqref{equ:grpo} are set to 0.2 and 0.27, correspondingly—these values have been precisely tuned for a fair assessment. It is important to highlight that GRPO requires the Routing Replay strategy to achieve stable convergence for MoE RL, as will be elaborated in \S~\ref{subsec:routing_replay}; in contrast, \textit{GSPO eliminates the necessity for such a strategy}.

As illustrated in Figure~\ref{fig:results}, the GSPO training process remains stable across all phases. Notably, \textbf{GSPO enables ongoing improvements in performance by scaling up training computation, systematically refreshing the query set, and increasing the output sequence length}. In addition, GSPO exhibits enhanced training efficiency compared to GRPO, attaining higher accuracy and stronger benchmark outcomes using equivalent amounts of training computation and query resources. Lastly, deploying GSPO in RL training for the latest Qwen3 models further validates GSPO’s effectiveness in realizing the benefits of RL scaling for large language models.

\subsection{Curious Observation on Clipping Fractions}

An important difference between GSPO and GRPO lies in GSPO's approach of applying clipping to complete sequences rather than on a per-token basis. As illustrated in Figure~\ref{fig:clipping}, there is a substantial gap—approximately two orders of magnitude—between the proportion of clipped tokens in GSPO and GRPO, and modifying the clipping thresholds does not mitigate this large discrepancy. Interestingly, even though GSPO clips a much greater share of tokens and consequently relies on a reduced set of tokens for both training and gradient estimation, it nonetheless attains superior training efficiency compared to GRPO. This somewhat paradoxical result—where more extensive clipping results in better training outcomes—suggests that GRPO’s method of estimating gradients at the token level introduces considerable noise and hampers sample efficiency. On the other hand, the sequence-level strategy of GSPO generates a more consistent and productive learning signal.

\subsection{Benefit of GSPO for MoE Training}
\label{subsec:routing_replay}

\paragraph{Background}
Relative to the RL training of dense neural networks, the inherently sparse activation characteristic of MoE models brings about distinct stability issues. Specifically, our observations indicate that under the GRPO algorithm, MoE models exhibit substantial \textit{expert-activation volatility}, which can impede the convergence of RL training. More precisely, we noticed that following one or more gradient steps, the selection of experts triggered by the same input often shifts significantly. For instance, in the 48-layer Qwen3-30B-A3B-Base model, approximately 10\% of the experts selected for a given rollout sample change after every RL gradient update, comparing the newly learned policy $\pi_{\theta}$ with the preceding policy $\pi_{\theta_\text{old}}$. This behavior, particularly evident in deeper MoE models, causes the token-wise importance weights $w_{i,t}(\theta) = \frac{ \pi_{\theta} (y_{i,t} | x, y_{i,<t}) }{ \pi_{\theta_\text{old}} (y_{i,t} | x,y_{i,<t})}$ to become highly unstable, undermining their reliability as described in \S~\ref{sec:motivation} and~\ref{subsec:gradient}, and consequently disrupts the proper convergence of RL training.

\paragraph{Our Previous Approach} In earlier work, we addressed this issue by utilizing the \textbf{Routing Replay} method during training. The core idea involved storing the set of experts activated under $\pi_{\theta_\text{old}}$ and subsequently reapplying these same routing selections within $\pi_{\theta}$ as we calculated the importance weights $w_{i,t}(\theta) = \frac{ \pi_{\theta} (y_{i,t} | x, y_{i,<t}) }{ \pi_{\theta_\text{old}} (y_{i,t} | x,y_{i,<t})}$. This procedure ensures that, for every token $y_{i,t}$, the probability computations for both $\pi_{\theta}$ and $\pi_{\theta_\text{old}}$ utilize the identical expert subset. Consequently, this alignment mitigates volatility in the computed token-level importance ratios and consistently trains the same underlying subnetwork through successive gradient steps. Figure~\ref{fig:routing_replay} illustrates that Routing Replay is pivotal for stabilizing and supporting the standard convergence behavior observed during GRPO training for MoE architectures.

\paragraph{Benefit of GSPO} While Routing Replay enables GRPO training for MoE models to achieve convergence, its reliance on storing and reusing routing decisions leads to greater demands on memory and inter-device communication, which may restrict the usable capacity of the MoE model. Alternatively, as illustrated in Figure~\ref{fig:results}, GSPO removes the requirement for Routing Replay and independently computes importance ratios $s_i(\theta)$ in the standard way, ensuring typical convergence behavior and consistent optimization. The central observation is that GSPO considers only the sequence likelihood (i.e., $\pi_{\theta} (y_i | x)$), rendering it insensitive to the likelihood of individual tokens (i.e., $\pi_{\theta} (y_{i,t} | x, y_{i,<t})$). Given that the MoE model retains stable language modeling ability throughout, sequence likelihood values remain relatively steady. Thus, GSPO offers a foundational solution to the instability of expert activation in MoE architectures, eliminating the necessity for intricate methods such as Routing Replay. This approach streamlines and enhances the robustness of the training process while enabling the model to fully utilize its inherent capacity, free from imposed limitations.

\subsection{Benefit of GSPO for RL Infrastructure}

Due to the differences in precision between training engines such as Megatron and inference engines like SGLang and vLLM, it is standard practice to recompute the likelihoods of sampled responses under the previous policy $\pi_{\theta_\text{old}}$ using the training engine. Nevertheless, GSPO employs optimization based on sequence-level likelihoods rather than token-level ones. Sequence-level likelihoods, in general, exhibit greater robustness to precision mismatches. Consequently, GSPO enables the direct utilization of likelihoods produced by the inference engine during optimization, eliminating the necessity for recomputation by the training engine. This approach is particularly advantageous in contexts such as partial rollout, multi-turn RL, and in frameworks where training and inference are separated.

\section{Conclusion}

We introduce Group Sequence Policy Optimization (GSPO), a novel reinforcement learning method tailored for large language model training. Adhering to the core concept of importance sampling, GSPO constructs importance ratios grounded in sequence probabilities and executes clipping, reward assignment, and optimization at the sequence level. Compared with GRPO, GSPO achieves significantly improved stability, sample efficiency, and overall performance during training. It proves especially effective for large-scale RL applications involving MoE models and forms the technological basis for the notable gains seen in recent Qwen3 model iterations. As a robust and scalable solution, GSPO paves the way for further expansion in RL, supporting future breakthroughs in machine intelligence.

\end{document}